\section*{Hoofdstuk \ref{ch:g10:universal-dna} --- DNA: een universeel erfelijk molecuul}

\begin{solution}{exo:g10:universal-dna:1}
Een fosfaatgroep, een suiker (desoxyribose) en een base. Basen: adenine,
thymine, guanine, cytosine.
\end{solution}

\begin{solution}{exo:g10:universal-dna:2}
\texttt{ATGCCTAAGT}, base onder base geschreven.
\end{solution}

\begin{solution}{exo:g10:universal-dna:3}
C $= 21\%$; A $=$ T $= (100 - 42)/2 = 29\%$ elk.
\end{solution}

\begin{solution}{exo:g10:universal-dna:4}
Een \omterm{def:g10:universal-dna:gene}{gen} is een stuk \omterm{def:g10:universal-dna:information}{DNA} waarvan de volgorde de instructies voor één
\omterm{def:g10:chemistry-of-life:families}{eiwit} draagt, op een vaste plaats op een \omterm{def:g10:universal-dna:information}{chromosoom}. Een \omterm{def:g10:universal-dna:gene}{allel} is een
van de versies van een \omterm{def:g10:universal-dna:gene}{gen}, die in enkele basen van haar volgorde van de
andere verschilt.
\end{solution}

\begin{solution}{exo:g10:universal-dna:5}
Een organisme in het \omterm{def:g10:universal-dna:information}{DNA} waarvan een \omterm{def:g10:universal-dna:gene}{gen} van een andere soort is
ingebracht, en dat het leest en doorgeeft. Voorbeelden: muizen met het
\omterm{def:g10:universal-dna:gene}{GFP-gen} van de kwal; \emph{E.~coli} met het menselijke insulinegen.
\end{solution}

\begin{solution}{exo:g10:universal-dna:6}
$1500 \times 0.34 = \qty{510}{nm}$; $1500/10 = 150$ windingen.
\end{solution}

\begin{solution}{exo:g10:universal-dna:7}
$2 \times \num{3.2e9} \times \qty{0.34}{nm} \approx \qty{2.2}{m}$,
ongeveer $\num{3.6e5}$ keer de doorsnede van de \omterm{def:g10:cells-common-unit:organelle}{celkern}.
\end{solution}

\begin{solution}{exo:g10:universal-dna:8}
Een "gezuiverde" fractie kan nog sporen van andere moleculen bevatten,
zodat haar werking op zichzelf niet bewijst dat \omterm{def:g10:universal-dna:information}{DNA} verantwoordelijk is.
Specifiek het \omterm{def:g10:universal-dna:information}{DNA} afbreken in het volledige extract en zien dat de
werking verdwijnt, terwijl enzymen die \omterm{def:g10:chemistry-of-life:families}{eiwit} en RNA afbreken haar
onaangetast laten, legt het effect vast op \omterm{def:g10:universal-dna:information}{DNA} en op niets anders.
\end{solution}

\begin{solution}{exo:g10:universal-dna:9}
Een vaste herhaling zou in elke soort dezelfde samenstelling geven (25\%
van elke base). Chargaff vond dat A+T sterk verschilt tussen soorten: de
volgorde van de basen is geen vast patroon, dus zij is vrij om te
variëren --- en juist dat maakt haar geschikt om informatie te dragen.
\end{solution}

\begin{solution}{exo:g10:universal-dna:10}
De volgorde wordt op volgorde gelezen om het \omterm{def:g10:chemistry-of-life:families}{eiwit} te bouwen; een
veranderde base op een plaats die een \omterm{def:g10:chemistry-of-life:families}{aminozuur} aanwijst, kan het ene
\omterm{def:g10:chemistry-of-life:families}{aminozuur} door het andere vervangen en de vorm of de werking van het
\omterm{def:g10:chemistry-of-life:families}{eiwit} veranderen. Ligt de veranderde base in een deel van het \omterm{def:g10:universal-dna:gene}{gen} dat
niet wordt gelezen, of wijst zij hetzelfde \omterm{def:g10:chemistry-of-life:families}{aminozuur} aan, dan blijft het
\omterm{def:g10:chemistry-of-life:families}{eiwit} ongewijzigd.
\end{solution}

\begin{solution}{exo:g10:universal-dna:11}
Alle \omterm{def:g10:cells-common-unit:cell}{cellen} van de muis stammen door deling af van de bevruchte eicel;
bij elke deling wordt het \omterm{def:g10:universal-dna:information}{DNA}, het transgen inbegrepen, gekopieerd
(\omterm{prop:g10:universal-dna:helix}{complementariteit}) en verdeeld. Huid-, lever- en zaadcellen dragen het
\omterm{def:g10:universal-dna:gene}{gen} dus alle --- en de zaadcel geeft het aan de volgende generatie door.
\end{solution}

\begin{solution}{exo:g10:universal-dna:12}
$\num{4.6e6} \times \qty{0.34}{nm} \approx \qty{1.6}{mm}$, zo'n 800 keer
de lengte van de \omterm{def:g10:cells-common-unit:cell}{cel}. Maar het molecuul is slechts \qty{2}{nm} breed:
zijn volume, ongeveer $\pi \times 1^2 \times \num{1.6e6}\,
\unit{nm^3} \approx \num{5e-3}\,\unit{\micro m^3}$, is ongeveer 1\% van
dat van de \omterm{def:g10:cells-common-unit:cell}{cel}; opgevouwen en opgerold past het er ruim in.
\end{solution}

\begin{solution}{exo:g10:universal-dna:13}
Menselijk \omterm{def:g10:universal-dna:information}{DNA} is 700 keer langer maar draagt slechts ongeveer 5 keer
zoveel \omterm{def:g10:universal-dna:gene}{genen}: het grootste deel bestaat niet uit \omterm{def:g10:universal-dna:gene}{genen}. Een deel van het
extra \omterm{def:g10:universal-dna:information}{DNA} regelt wanneer \omterm{def:g10:universal-dna:gene}{genen} worden gebruikt, een deel is herhaalde
volgorde zonder bekende functie. Het aantal basen meet de omvang van de
tekst, niet de ingewikkeldheid van het organisme; het aantal \omterm{def:g10:universal-dna:gene}{genen} en
vooral de manier waarop ze worden gecombineerd en geregeld, wegen
zwaarder.
\end{solution}

\begin{solution}{exo:g10:universal-dna:14}
Het zou de universaliteit van het \omterm{def:g10:universal-dna:information}{DNA} tegenspreken
(\cref{prop:g10:universal-dna:universal}). Omdat een gedeeld molecuul en
een gedeelde code van een gemeenschappelijke voorouder zijn geërfd, zou
zo'n soort uit een afzonderlijke oorsprong van het leven moeten
voortkomen: haar \omterm{def:g10:universal-dna:gene}{genen} zouden niet door onze \omterm{def:g10:cells-common-unit:cell}{cellen} gelezen kunnen
worden en \omterm{prop:g10:universal-dna:universal}{transgenese} tussen de twee lijnen zou onmogelijk zijn.
\end{solution}

\begin{solution}{exo:g10:universal-dna:15}
Het bacteriële product is precies het menselijke \omterm{def:g10:chemistry-of-life:families}{eiwit}
(varkensinsuline verschilt in één \omterm{def:g10:chemistry-of-life:families}{aminozuur} en kan afweerreacties
uitlokken), en het wordt in onbeperkte hoeveelheid gemaakt, vrij van
varkensziekteverwekkers. Een toezichthouder zou nog altijd vragen of het
gezuiverde \omterm{def:g10:chemistry-of-life:families}{eiwit} vrij is van bacteriële verontreinigingen en of het
correct gevouwen en werkzaam is.
\end{solution}

\begin{solution}{pb:g10:universal-dna:1}
\textbf{1.} $720 \times 0.34 \approx \qty{245}{nm}$; $720/10 = 72$
windingen.

\textbf{2.} \texttt{TACTCATTTCCTCTTCTTG}.

\textbf{3.} T $= 27\%$; G $=$ C $= (100 - 54)/2 = 23\%$.

\textbf{4.} $720/238 \approx 3.0$: drie basen per \omterm{def:g10:chemistry-of-life:families}{aminozuur} (de laatste
drie geven het einde van de keten aan).

\textbf{5.} $4^{720}$ reeksen, ongeveer $10^{433}$. De chemie laat elke
volgorde toe, dus de ene volgorde die uit al die mogelijkheden gekozen
is, is een boodschap --- informatie.

\textbf{6.} De universaliteit: elk levend wezen gebruikt dezelfde vier
\omterm{def:g10:universal-dna:nucleotide}{nucleotiden} en leest een volgorde volgens dezelfde regels.

\textbf{7.} Alle \omterm{def:g10:cells-common-unit:cell}{cellen} van de muis stammen door deling van de eicel af,
dus een \omterm{def:g10:universal-dna:gene}{gen} dat in de eicel is ingebracht, wordt naar elke \omterm{def:g10:cells-common-unit:cell}{cel}
gekopieerd. In de huid van een volwassen dier ingebracht, zou het alleen
in de behandelde huidcellen en hun nakomelingen zitten.

\textbf{8.} Elke celsoort bevat het volledige \omterm{def:g10:universal-dna:information}{DNA}, het transgen
inbegrepen: specialisatie gooit geen \omterm{def:g10:universal-dna:gene}{genen} weg.

\textbf{9.} Ja: de geslachtscellen dragen het \omterm{def:g10:universal-dna:gene}{gen}. Heeft de muis één
kopie op één \omterm{def:g10:universal-dna:information}{chromosoom} van een paar, dan erft ongeveer de helft van
haar nakomelingen het en gloeit.

\textbf{10.} Het \omterm{def:g10:chemistry-of-life:families}{eiwit} is geen informatie: het wordt opgebruikt en nooit
vernieuwd, dus dooft de gloed. Het \omterm{def:g10:universal-dna:gene}{gen} wordt bij elke deling gekopieerd
en voortdurend gelezen, dus is de gloed blijvend en erfelijk.

\textbf{11.} De DNA-fractie; behandeling met een enzym dat \omterm{def:g10:universal-dna:information}{DNA} afbreekt.

\textbf{12.} Bacteriën namen \omterm{def:g10:universal-dna:information}{DNA} van een andere stam op, bouwden het in,
brachten het tot uiting en gaven het kenmerk door: een \omterm{def:g10:universal-dna:gene}{gen} dat van het
ene organisme naar het andere is overgebracht, en dat is precies wat
\omterm{prop:g10:universal-dna:universal}{transgenese} doelbewust doet.

\textbf{13.} A$\,=\,$T en G$\,=\,$C ondersteunen de paring; het
wisselende aandeel A+T laat zien dat de volgorde van de paren vrij is.

\textbf{14.} A en G zijn grote basen met twee ringen, T en C kleine met
één ring: een grote gepaard met een kleine overspant altijd dezelfde
afstand, zodat de twee ruggengraten over het hele molecuul \qty{2}{nm}
uit elkaar blijven. Twee grote basen zouden uitpuilen, twee kleine een
gat laten.

\textbf{15.} Omdat elke streng het complement van de andere is: scheid
ze, bouw op elk volgens de paringsregel een partner, en er ontstaan twee
gelijke moleculen.

\textbf{16.} $2 \times \num{2.7e9} \times \qty{0.34}{nm} \approx
\qty{1.8}{m}$.

\textbf{17.} $\qty{1.8}{m} / \qty{6}{\micro m} = \num{3e5}$: een
driehonderdduizendvoudige samenpakking.

\textbf{18.} $720 / \num{5.4e9} \approx \num{1.3e-7}$: ongeveer een
tienmiljoenste van het \omterm{def:g10:universal-dna:information}{DNA} van de \omterm{def:g10:cells-common-unit:cell}{cel}.

\textbf{19.} $\num{2e4} \times \num{3e4} = \num{6e8}$ basenparen,
ongeveer 22\% van het genoom; de rest ligt tussen de \omterm{def:g10:universal-dna:gene}{genen}.

\textbf{20.} Een muizencel draagt ongeveer \qty{1.8}{m} \omterm{def:g10:universal-dna:information}{DNA}; omdat dat
\omterm{def:g10:universal-dna:information}{DNA} precies zo gebouwd is en gelezen wordt als dat van de kwal, wordt
het ene vreemde \omterm{def:g10:universal-dna:gene}{gen} tussen de twintigduizend andere net als de rest
gelezen.
\end{solution}
